\begin{afterword}
\summaryhead

The post imagines physicists meeting quantum mechanics before what it calls the correct answer, many-worlds, was available. The best they could have done, it says, was a strict ``shut up and calculate'': affirm only what the data and calculations force, so that no later theory could contradict them.

Instead, in dialogues the post says may never have happened, the careful Gallant's ``I don't know what these equations mean'' becomes Goofus's ``these equations don't mean anything,'' and a rule for calculating after a measurement becomes the claim that knowledge makes amplitudes vanish. Read literally, that is consciousness causing collapse, which the post finds implausible but at least a picture of reality. The anti-realist who calls the equations meaningless, and who answers questions about entangled photons with ``It's meaningless to talk about it,'' is asked how he knows. The post argues that ``meaningless'' stops inquiry and that the equations' success suggests they mean something. Its moral: it is hard to stay openly confused.

\responsehead

The post's question is a fair one. Someone who says the Schrödinger equation ``doesn't mean anything'' owes an account of how they know, and the post limits its scorn to those who turned ``We don't know why it works'' into that claim. Its main argument, that the equations' success is best explained by their meaning something, is the ``no miracles'' argument for scientific realism, in the form Putnam gave it in 1975. It is a respectable argument with well-known replies, which the post does not mention.

The opponents, though, are left unnamed, and the post says its dialogues may never have happened. The views they mock came with reasons. Bohr treated the wave function as a tool for prediction because it lives in configuration space and involves imaginary numbers. Pauli compared asking whether something exists that one cannot know anything about to asking how many angels can sit on a needle. QBism, developed from 2002, treats a quantum state as an agent's degrees of belief and entangled correlations as updates of expectations, with no nonlocal influence. All three are disputed. But each answers the post's two questions, ``Where is Goofus getting all this stuff?'' and what reality would make his words true, and the post considers none of them. The chain of reasoning it calls ``one of the most embarrassing wrong turns in the history of science'' is close to QBism, whose reading of probability resembles the post's own.

The post's certainty conflicts with its own rule. It asks the strict calculator to affirm nothing ``not forced by the data and the calculations,'' says of decoherence only that it ``is not ruled out,'' and allows that the correct theory may be something else. It also calls many-worlds ``the correct answer'' and the source of ``a certainty of many worlds existing.'' A strict calculator would not make that claim, and the post's grounds for it are not in this post. The full case for many-worlds is argued later in the sequence; here it is assumed.

Smaller points. As far as the record shows, Einstein put the epigraph's question to Abraham Pais, not to Bohr. Bell's theorem rules out local hidden variables, not deterministic ones: Bohm's theory is consistent with it as a nonlocal theory. The claim about textbooks has no example; the claim that almost nobody admits to believing that consciousness causes collapse is supported by a later poll.

In short: a fair question for anyone who calls the equations meaningless, put to opponents the post does not name, without the reasons the actual anti-realists gave, by a post that calls many-worlds ``a certainty'' beside a rule against affirming what the data do not force.
\end{afterword}
